\section{Making Useful Measurements}\label{sec:measurements}

\paragraph{Establish a reference}
For a controlled check at $48\,\mathrm{kHz}$, use a $750\,\mathrm{Hz}$ sine:
with Length 2048 it falls exactly on bin 32. Use a known $5\,\mathrm{V}$
peak ($10\,\mathrm{V}$ peak-to-peak) source, unity input gain, Flattop,
zero Slope, no smoothing, and AC coupling off. After the frame fills,
the peak should be approximately $0\,\mathrm{dB}$ or $100\,\%$.
Use the source's frequency setting or a separate tuner to establish pitch;
the cursor does not snap to a peak. An off-bin tone can read differently.

\begin{center}
\renewcommand{\arraystretch}{1.25}
\begin{tabular}{@{}rrr@{}}
\toprule
\textbf{Sine peak voltage} & \textbf{Relative amplitude} & \textbf{Level} \\
\midrule
$5\,\mathrm{V}$ & $100\,\%$ & $0\,\mathrm{dB}$ \\
$2.5\,\mathrm{V}$ & $50\,\%$ & $-6.02\,\mathrm{dB}$ \\
$0.5\,\mathrm{V}$ & $10\,\%$ & $-20\,\mathrm{dB}$ \\
$0.005\,\mathrm{V}$ & $0.1\,\%$ & $-60\,\mathrm{dB}$ \\
\bottomrule
\end{tabular}
\end{center}

These ideal isolated-tone values follow $20\log_{10}(V_{\rm peak}/5\,\mathrm{V})$.
They are not a conversion from an arbitrary signal's largest spectral peak
to its total voltage. This reference is not dBFS, SPL, LUFS, or a calibrated
noise-power density measurement.

\paragraph{Know what is being summed}
Each jack sums all voices of a polyphonic cable before analysis. Two identical,
in-phase voices add to twice the amplitude (about $+6.02\,\mathrm{dB}$);
equal opposite-polarity voices cancel. To inspect voices independently, split
them into separate monophonic cables and use different inputs. Equal-looking
magnitude spectra alone cannot tell you the relative phase or predict the
result of mixing two signals.

\paragraph{Separate signal changes from viewing changes}
Input gain and AC coupling affect newly captured samples. Window, Average,
and Smooth affect the analysis. Frequency bounds and scales determine the
view; Slope deliberately weights levels by frequency. Turn Run off to compare
analysis choices on retained audio, but remember that Fourier keeps
reanalyzing and Average can still settle. Changing Length clears that audio.

\paragraph{Record the conditions}
When comparing a patch revision, note the Rack sample rate, input source and
gain, Length, Hop, window, smoothing, Slope, scales, and AC coupling. A noise
trace can change height when Length or window changes because each bin covers
a different effective bandwidth. Do not diagnose a quieter source from that
change alone. Save settings as a preset; save an external image if you need a
record of the trace, because patches do not retain captured spectra.
